\documentclass[tikz,border=20pt]{standalone}
\usetikzlibrary{calc, arrows.meta, decorations.pathmorphing}
\usepackage[UTF8]{ctex}

\begin{document}
\begin{tikzpicture}[x=1cm, y=1cm,
   iso/.style={very thick, draw=green!45!black},
]

% =====================================================================
% 纯 2D 线框等值面 demo：保持平面轮廓风格，但拓扑更丰富
% =====================================================================

% ---------------------------------------------------------------
% 宏 1：单洞环面 g=1（双椭圆嵌套，原风格）
% ---------------------------------------------------------------
\newcommand{\torusone}[3]{%
  \begin{scope}[shift={(#1,#2)}, scale=#3]
    \draw[iso] (0,0) ellipse (0.85 and 0.55);
    \draw[iso] (0,0) ellipse (0.32 and 0.20);
  \end{scope}
}

% ---------------------------------------------------------------
% 宏 2：双洞曲面 g=2（外轮廓 + 两个洞）
% ---------------------------------------------------------------
\newcommand{\torustwo}[3]{%
  \begin{scope}[shift={(#1,#2)}, scale=#3]
    \draw[iso] (0,0) ellipse (1.15 and 0.55);
    \draw[iso] (-0.5,0) ellipse (0.26 and 0.18);
    \draw[iso] ( 0.5,0) ellipse (0.26 and 0.18);
  \end{scope}
}

% ---------------------------------------------------------------
% 宏 3：三洞曲面 g=3
% ---------------------------------------------------------------
\newcommand{\torusthree}[3]{%
  \begin{scope}[shift={(#1,#2)}, scale=#3]
    \draw[iso] (0,0) ellipse (1.45 and 0.55);
    \draw[iso] (-0.8,0) ellipse (0.24 and 0.17);
    \draw[iso] ( 0.0,0) ellipse (0.24 and 0.17);
    \draw[iso] ( 0.8,0) ellipse (0.24 and 0.17);
  \end{scope}
}

% ---------------------------------------------------------------
% 宏 4：光滑闭合曲线 g=0（单一轮廓，无洞——最简单球面截线）
% ---------------------------------------------------------------
\newcommand{\loopzero}[3]{%
  \begin{scope}[shift={(#1,#2)}, scale=#3]
    \draw[iso] (0,0) circle (0.62);
  \end{scope}
}

% ---------------------------------------------------------------
% 宏 5：不规则闭合轮廓 g=0（复杂几何、带凹凸的等值线）
% ---------------------------------------------------------------
\newcommand{\blobzero}[3]{%
  \begin{scope}[shift={(#1,#2)}, scale=#3]
    \draw[iso] plot[smooth cycle, tension=0.85]
      coordinates {(0,0.68)(0.5,0.52)(0.42,0.05)(0.75,-0.32)
                   (0.28,-0.6)(-0.2,-0.42)(-0.6,-0.55)(-0.7,0.0)
                   (-0.42,0.35)(-0.55,0.62)};
  \end{scope}
}

% ---------------------------------------------------------------
% 宏 6：多分量等值线 c=3（三条独立闭合曲线，连通性示意）
% ---------------------------------------------------------------
\newcommand{\multicomp}[3]{%
  \begin{scope}[shift={(#1,#2)}, scale=#3]
    \draw[iso] (-0.55,0.25) circle (0.34);
    \draw[iso] ( 0.5,0.4) circle (0.30);
    \draw[iso] ( 0.05,-0.45) circle (0.38);
  \end{scope}
}

% ===== 摆放展示 =====
\loopzero{0}{0}{1}
\node[font=\sffamily\footnotesize, align=center, text=black!70] at (0,-1.4)
   {单一闭曲线\\$g=0$};

\torusone{2.4}{0}{1}
\node[font=\sffamily\footnotesize, align=center, text=black!70] at (2.4,-1.4)
   {单洞环面\\$g=1$};

\torustwo{5.0}{0}{1}
\node[font=\sffamily\footnotesize, align=center, text=black!70] at (5.0,-1.4)
   {双洞曲面\\$g=2$};

\torusthree{8.0}{0}{1}
\node[font=\sffamily\footnotesize, align=center, text=black!70] at (8.0,-1.4)
   {三洞曲面\\$g=3$};

\blobzero{11.0}{0}{1}
\node[font=\sffamily\footnotesize, align=center, text=black!70] at (11.0,-1.4)
   {不规则轮廓\\$g=0$（复杂几何）};

\multicomp{13.6}{0}{1}
\node[font=\sffamily\footnotesize, align=center, text=black!70] at (13.6,-1.4)
   {多分量\\$c=3$};

\node[font=\sffamily\bfseries\large] at (6.8,2.2)
   {纯 2D 线框等值面 demo（保持平面风格，拓扑更丰富）};

\end{tikzpicture}
\end{document}